%% kapa: exam=2 round=21 year=2010 period=1 subject=appraisal_practice subject_ko=감정평가실무 seq=1 number=1 points=40
%% source: Q-Net 감정평가사 기출문제 (cst00309) artlSeq=5024183 file=extracted/1008/5024183/2010년 제21회 감정평가사 2차시험 기출문제/1교시.hwp
%% conversion: kordoc 4.24.51 -> tools/gen2.py (auto); review=proofread-sonnet; flags=figure,tikz
\begin{longquestion}{【문제 1】}
공정감정평가법인 소속 감정평가사인 김한국씨는 아래 부동산중 이대한씨 지분에 대해서 한강은행(담보)과 甲구청(보상)으로부터 동시에 평가의뢰를 받고 감정평가가격을 구한 후 감정평가서를 작성하고자 한다. 주어진 자료를 활용하여 다음의 물음에 답하시오. (40점) \par
물음(1) 평가목적별로 이대한씨 소유의 토지와 건물을 평가하시오. (20점) \par
1) 평가목적이 담보일 경우 \par
2) 평가목적이 보상일 경우 \par
물음(2) 감정평가에 관한 규칙 제9조에서 규정하고 있는 필수적 기재사항에 의거 서술식으로 감정평가서를 작성하시오(단, 평가목적이 보상일 경우에 중복되는 항목은 생략). (20점) \par
1) 평가목적이 담보일 경우 \par
2) 평가목적이 보상일 경우 \par
<자료 1> 사전조사사항 \par
1) 토지등기부등본 \par
{\small\centering
\begin{tabular}{|c|c|c|c|c|}
\hline
\makecell{기호} & \makecell{소재지} & \makecell{지목} & \makecell{면적(m\textsuperscript{2})} & \makecell{소유자} \\
\hline
\makecell{1} & \makecell{甲구 乙동 54} & \makecell{대} & \makecell{500} & \makecell{공유자지분 2분의 1 이대한\newline{}공유자지분 2분의 1 박조선} \\
\hline
\makecell{2} & \makecell{甲구 乙동 산75} & \makecell{임야} & \makecell{2,550} & \makecell{공유자지분 2분의 1 이대한\newline{}공유자지분 2분의 1 박조선} \\
\hline
\end{tabular}\par}

2) 건물등기부등본 \par
{\small\centering
\begin{tabular}{|c|c|c|c|c|c|}
\hline
\makecell{기호} & \makecell{소재지} & \makecell{물건의 종류} & \makecell{구조, 규격} & \makecell{면적(m\textsuperscript{2})} & \makecell{소유자} \\
\hline
\makecell{가} & \makecell{甲구 乙동 54} & \makecell{점포} & \makecell{철근콘크리트조 슬라브지붕 단층} & \makecell{80} & \makecell{이대한} \\
\hline
\makecell{나} & \makecell{甲구 乙동 54} & \makecell{주택} & \makecell{시멘트벽돌조 슬라브지붕 단층} & \makecell{70} & \makecell{박조선} \\
\hline
\end{tabular}\par}

3) 토지대장 \par
{\small\centering
\begin{tabular}{|c|c|c|c|c|}
\hline
\makecell{기호} & \makecell{소재지} & \makecell{지목} & \makecell{면적(m\textsuperscript{2})} & \makecell{소유자} \\
\hline
\makecell{1} & \makecell{甲구 乙동 54} & \makecell{대} & \makecell{500} & \makecell{공유자지분 2분의 1 이대한\newline{}공유자지분 2분의 1 박조선} \\
\hline
\makecell{2} & \makecell{甲구 乙동 산75} & \makecell{임야} & \makecell{2,800} & \makecell{공유자지분 2분의 1 이대한\newline{}공유자지분 2분의 1 박조선} \\
\hline
\end{tabular}\par}

4) 건축물대장 \par
{\small\centering
\begin{tabular}{|c|c|c|c|c|c|}
\hline
\makecell{기호} & \makecell{소재지} & \makecell{물건의 종류} & \makecell{구조, 규격} & \makecell{면적(m\textsuperscript{2})} & \makecell{소유자} \\
\hline
\makecell{가} & \makecell{甲구 乙동 54} & \makecell{점포} & \makecell{철근콘크리트조 슬라브지붕 단층} & \makecell{80} & \makecell{이대한} \\
\hline
\makecell{나} & \makecell{甲구 乙동 54} & \makecell{주택} & \makecell{시멘트벽돌조 슬라브지붕 단층} & \makecell{70} & \makecell{박조선} \\
\hline
\end{tabular}\par}

5) 甲구청 제시목록 \par
① 토지 \par
{\small\centering
\begin{tabular}{|c|c|c|c|c|}
\hline
\makecell{기호} & \makecell{소재지} & \makecell{편입면적(m\textsuperscript{2})} & \makecell{지목} & \makecell{비고} \\
\hline
\makecell{1} & \makecell{甲구 乙동 54} & \makecell{100} & \makecell{대} & \makecell{소유자 : 이대한} \\
\hline
\makecell{2} & \makecell{甲구 乙동 산75} & \makecell{500} & \makecell{임야} & \makecell{공유자지분 2분의 1 이대한\newline{}공유자지분 2분의 1 박조선} \\
\hline
\end{tabular}\par}

② 건물 \par
{\small\centering
\begin{tabular}{|c|c|c|c|c|c|}
\hline
\makecell{기호} & \makecell{소재지} & \makecell{물건의 종류} & \makecell{구조, 규격} & \makecell{편입면적(m\textsuperscript{2})} & \makecell{비고} \\
\hline
\makecell{가} & \makecell{甲구 乙동 54} & \makecell{점포} & \makecell{철근콘크리트조\newline{}슬라브지붕 단층} & \makecell{20} & \makecell{소유자 : 이대한} \\
\hline
\makecell{㉠} & \makecell{甲구 乙동 54} & \makecell{창고} & \makecell{시멘트블럭조\newline{}슬라브지붕 단층} & \makecell{10} & \makecell{소유자 : 이대한} \\
\hline
\end{tabular}\par}

6) 토지이용계획확인원 \par
① 기호1 : 일반상업지역(250m\textsuperscript{2}), 제2종일반주거지역(250m\textsuperscript{2}), 도시계획도로저촉 \par
② 기호2 : 자연녹지지역 \par
7) 지적도 \par
\qfigure{figures/01-1.png}
<자료 2> 현장조사사항 \par
1) 2010. 08. 30에 현장조사를 하였으나 가격자료수집이 미흡하여 2010. 09. 02에 재조사 완료하였음. \par
2) 기호1 토지는 인접필지와 대체로 평탄하며, 점포부지(이대한 소유)와 주택부지(박조선 소유)로 이용 중이고, 기호2 토지는 남하향의 완경사 자연림으로서 형상은 부정형이고, 지상에 자연생 활잡목이 자생하고 있으나 경제적 가치는 없는 것으로 판단됨. \par
3) 기호1 토지의 도로접면은 지적도와 동일하며, 기호2 토지의 도로접면은 세로가에 접하고 있음. \par
4) 기호1 토지의 지상에 시멘트블럭조 슬라브지붕 단층 창고(이대한 소유 : 기호㉠)와 시멘트블럭조 슬라브지붕 단층 창고(박조선 소유 : 기호㉡)가 무허가건축물로 존재하고 있으며, 신축년도는 모두 2001. 09. 07에 건축된 것으로 탐문조사 되었고, 이로 인하여 토지에 미치는 영향은 없는 것으로 판단됨. \par
5) 건물배치도 \par
\begin{center}
\begin{tikzpicture}[x=0.0125cm,y=-0.0125cm,line width=0.7pt,font=\small]
\draw (197,5)--(197,465);
\draw (313,3)--(313,465);
\draw (664,37)--(664,480);
\draw (692,40)--(692,480);
\draw[dashed] (360,12)--(360,435);
\foreach \y in {115,218,325,432} \draw (318,\y)--(664,\y);
\draw (488,115)--(488,218);
\node at (415,168) {53}; \node at (568,168) {53-1};
\node at (482,272) {54}; \node at (487,378) {55};
\draw (346,260) rectangle (390,320); \node at (369,288) {가};
\node[anchor=west] at (400,305) {80m\textsuperscript{2}};
\draw (542,273) rectangle (586,305); \node at (564,289) {나};
\node[anchor=west] at (594,300) {70m\textsuperscript{2}};
\draw[dashed] (318,228) rectangle (350,252); \node[font=\tiny] at (334,240) {㉠};
\draw[dashed] (500,230) rectangle (530,252); \node[font=\tiny] at (515,241) {㉡};
\node[align=center] at (110,282) {도시계획도로\\ 저촉부분 40m\textsuperscript{2}};
\draw[-{Stealth[length=2mm]}] (175,290)--(340,290);
\draw[<->] (313,458)--(360,458);
\node at (340,493) {도시계획도로};
\end{tikzpicture}
\end{center}
<자료 3> 가격결정을 위한 참고자료 \par
1) 표준지공시지가 현황 \par
{\small\centering
\setlength{\tabcolsep}{3pt}
\begin{tabular}{|c|c|c|c|c|c|c|c|c|c|c|}
\hline
\multirow{2}{*}{\makecell{기호}} & \multirow{2}{*}{\makecell{소재지}} & \multirow{2}{*}{\makecell{면적\newline{}(m\textsuperscript{2})}} & \multirow{2}{*}{\makecell{지목}} & \multirow{2}{*}{\makecell{이용\newline{}상황}} & \multirow{2}{*}{\makecell{용도지역}} & \multirow{2}{*}{\makecell{도로교통}} & \multirow{2}{*}{\makecell{형상\newline{}지세}} & \multicolumn{3}{|c|}{\makecell{공시지가(원/m\textsuperscript{2})}} \\
\cline{9-11}
 &  &  &  &  &  &  &  & \makecell{2008년} & \makecell{2009년} & \makecell{2010년} \\
\hline
\makecell{A} & \makecell{乙동 57} & \makecell{250} & \makecell{대} & \makecell{상업용} & \makecell{일반상업} & \makecell{광대세각} & \makecell{세장형 평지} & \makecell{900,000} & \makecell{980,000} & \makecell{1,100,000} \\
\hline
\makecell{B} & \makecell{乙동 58-1} & \makecell{250} & \makecell{대} & \makecell{단독\newline{}주택} & \makecell{제2종\newline{}일반주거} & \makecell{세로가} & \makecell{세장형평지} & \makecell{380,000} & \makecell{400,000} & \makecell{420,000} \\
\hline
\makecell{C} & \makecell{乙동\newline{}산74} & \makecell{3,000} & \makecell{임야} & \makecell{자연림} & \makecell{자연녹지} & \makecell{맹지} & \makecell{부정형 완경사} & \makecell{38,000} & \makecell{45,000} & \makecell{50,000} \\
\hline
\end{tabular}\par}

※ 표준지 공시지가는 공법상 제한이 없는 상태임. \par
2) 지가변동률(甲구) \par
{\small\centering
\begin{tabular}{|c|c|c|c|}
\hline
\makecell{구분} & \makecell{상업지역} & \makecell{주거지역} & \makecell{녹지지역} \\
\hline
\makecell{2008. 01. 01 - 2010. 08. 30} & \makecell{1.15100} & \makecell{1.20000} & \makecell{1.25000} \\
\hline
\makecell{2008. 01. 01 - 2010. 09. 02} & \makecell{1.15500} & \makecell{1.21000} & \makecell{1.25500} \\
\hline
\makecell{2009. 01. 01 - 2010. 08. 30} & \makecell{1.10000} & \makecell{1.12000} & \makecell{1.13000} \\
\hline
\makecell{2009. 01. 01 - 2010. 09. 02} & \makecell{1.11500} & \makecell{1.12500} & \makecell{1.13500} \\
\hline
\makecell{2010. 01. 01 - 2010. 08. 30} & \makecell{1.05000} & \makecell{1.06500} & \makecell{1.07000} \\
\hline
\makecell{2010. 01. 01 - 2010. 09. 02} & \makecell{1.05500} & \makecell{1.07000} & \makecell{1.07500} \\
\hline
\end{tabular}\par}

3) 지역요인 : 동일함. \par
4) 개별요인 \par
① 도로접면 \par
{\small\centering
\begin{tabular}{|c|c|c|c|c|c|c|}
\hline
\makecell{구 분} & \makecell{광대세각} & \makecell{광대한면} & \makecell{소로한면} & \makecell{세로가} & \makecell{세로(불)} & \makecell{맹지} \\
\hline
\makecell{광대세각} & \makecell{1.00} & \makecell{0.95} & \makecell{0.86} & \makecell{0.81} & \makecell{0.75} & \makecell{0.72} \\
\hline
\makecell{광대한면} & \makecell{1.05} & \makecell{1.00} & \makecell{0.91} & \makecell{0.85} & \makecell{0.78} & \makecell{0.75} \\
\hline
\makecell{소로한면} & \makecell{1.16} & \makecell{1.10} & \makecell{1.00} & \makecell{0.93} & \makecell{0.86} & \makecell{0.83} \\
\hline
\makecell{세로가} & \makecell{1.24} & \makecell{1.18} & \makecell{1.07} & \makecell{1.00} & \makecell{0.92} & \makecell{0.89} \\
\hline
\makecell{세로(불)} & \makecell{1.34} & \makecell{1.28} & \makecell{1.16} & \makecell{1.09} & \makecell{1.00} & \makecell{0.96} \\
\hline
\makecell{맹지} & \makecell{1.39} & \makecell{1.32} & \makecell{1.20} & \makecell{1.12} & \makecell{1.04} & \makecell{1.00} \\
\hline
\end{tabular}\par}

② 형상 \par
{\small\centering
\begin{tabular}{|c|c|c|c|c|}
\hline
\makecell{구 분} & \makecell{정방형} & \makecell{장방형} & \makecell{사다리형} & \makecell{부정형} \\
\hline
\makecell{정방형} & \makecell{1.00} & \makecell{0.99} & \makecell{0.98} & \makecell{0.95} \\
\hline
\makecell{장방형} & \makecell{1.01} & \makecell{1.00} & \makecell{0.99} & \makecell{0.96} \\
\hline
\makecell{사다리형} & \makecell{1.02} & \makecell{1.01} & \makecell{1.00} & \makecell{0.97} \\
\hline
\makecell{부정형} & \makecell{1.05} & \makecell{1.04} & \makecell{1.03} & \makecell{1.00} \\
\hline
\end{tabular}\par}

③ 지세 \par
{\small\centering
\begin{tabular}{|c|c|c|}
\hline
\makecell{구 분} & \makecell{평지} & \makecell{완경사} \\
\hline
\makecell{평지} & \makecell{1.00} & \makecell{0.97} \\
\hline
\makecell{완경사} & \makecell{1.03} & \makecell{1.00} \\
\hline
\end{tabular}\par}

④ 공법상 제한(도시계획도로 저촉 : 토지, 건물 공동사용) \par
{\small\centering
\begin{tabular}{|c|c|c|}
\hline
\makecell{구 분} & \makecell{일반} & \makecell{제한} \\
\hline
\makecell{일반} & \makecell{1.00} & \makecell{0.85} \\
\hline
\makecell{제한} & \makecell{1.18} & \makecell{1.00} \\
\hline
\end{tabular}\par}

5) 기타요인 산정을 위한 자료 \par
① 평가선례 \par
{\small\centering
\begin{tabular}{|c|c|c|c|c|c|c|c|c|}
\hline
\makecell{기호} & \makecell{소재지} & \makecell{지목} & \makecell{면적\newline{}(m\textsuperscript{2})} & \makecell{이용상황} & \makecell{용도\newline{}지역} & \makecell{단가(원/m\textsuperscript{2})} & \makecell{가격시점} & \makecell{평가목적} \\
\hline
\makecell{A} & \makecell{乙동 59} & \makecell{대} & \makecell{250} & \makecell{상업용} & \makecell{일반\newline{}상업} & \makecell{1,300,000} & \makecell{2010. 01. 01} & \makecell{담보} \\
\hline
\makecell{B} & \makecell{乙동\newline{}59-1} & \makecell{대} & \makecell{250} & \makecell{단독주택} & \makecell{2종일반주거} & \makecell{500,000} & \makecell{2010. 01. 01} & \makecell{담보} \\
\hline
\makecell{C} & \makecell{乙동\newline{}60} & \makecell{대} & \makecell{250} & \makecell{상업용} & \makecell{일반\newline{}상업} & \makecell{1,150,000} & \makecell{2009. 01. 01} & \makecell{담보} \\
\hline
\makecell{D} & \makecell{乙동\newline{}60-1} & \makecell{대} & \makecell{250} & \makecell{단독주택} & \makecell{2종일반주거} & \makecell{480,000} & \makecell{2009. 01. 01} & \makecell{담보} \\
\hline
\makecell{E} & \makecell{乙동\newline{}110} & \makecell{대} & \makecell{250} & \makecell{상업용} & \makecell{일반\newline{}상업} & \makecell{1,500,000} & \makecell{2010. 01. 01} & \makecell{보상} \\
\hline
\makecell{F} & \makecell{乙동\newline{}110-1} & \makecell{대} & \makecell{250} & \makecell{단독주택} & \makecell{2종일반주거} & \makecell{600,000} & \makecell{2010. 01. 01} & \makecell{보상} \\
\hline
\makecell{G} & \makecell{乙동\newline{}111} & \makecell{대} & \makecell{250} & \makecell{상업용} & \makecell{일반\newline{}상업} & \makecell{1,300,000} & \makecell{2009. 01. 01} & \makecell{보상} \\
\hline
\makecell{H} & \makecell{乙동\newline{}111-1} & \makecell{대} & \makecell{250} & \makecell{단독주택} & \makecell{2종일반주거} & \makecell{560,000} & \makecell{2009. 01. 01} & \makecell{보상} \\
\hline
\makecell{I} & \makecell{乙동\newline{}112} & \makecell{대} & \makecell{250} & \makecell{상업용} & \makecell{일반\newline{}상업} & \makecell{1,100,000} & \makecell{2008. 01. 01} & \makecell{보상} \\
\hline
\makecell{J} & \makecell{乙동\newline{}112-1} & \makecell{대} & \makecell{250} & \makecell{단독주택} & \makecell{2종일반주거} & \makecell{540,000} & \makecell{2008. 01. 01} & \makecell{보상} \\
\hline
\end{tabular}\par}

② 지가변동률은 상기에서 제시한 자료와 동일하고 대상토지와 평가선례와의 지역요인 및 개별요인은 동일함. \par
③ 평가선례는 공법상 제한이 없는 상태임. \par
④ 지목이 ‘임야’인 표준지공시지가는 지목이 ‘대’인 표준지공시지가와 기타요인이 동일함. \par
<자료 4> 건물평가 자료 \par
1) 사용승인일자 \par
{\small\centering
\begin{tabular}{|c|c|c|}
\hline
\makecell{구분} & \makecell{점포} & \makecell{주택} \\
\hline
\makecell{사용승인일자} & \makecell{2001. 09. 07.} & \makecell{2001. 09. 07.} \\
\hline
\end{tabular}\par}

2) 건물 재조달원가(원/m\textsuperscript{2}) \par
{\small\centering
\begin{tabular}{|c|c|c|c|}
\hline
\makecell{구분} & \makecell{철근콘크리트조} & \makecell{시멘트벽돌조} & \makecell{시멘트블럭조} \\
\hline
\makecell{주택} & \makecell{1,000,000} & \makecell{900,000} & \makecell{600,000} \\
\hline
\makecell{점포} & \makecell{700,000} & \makecell{550,000} & \makecell{450,000} \\
\hline
\makecell{창고} & \makecell{500,000} & \makecell{400,000} & \makecell{350,000} \\
\hline
\end{tabular}\par}

3) 내용년수 \par
{\small\centering
\begin{tabular}{|c|c|c|c|}
\hline
\makecell{구분} & \makecell{철근콘크리트조} & \makecell{시멘트벽돌조} & \makecell{시멘트블럭조} \\
\hline
\makecell{내용년수} & \makecell{50년} & \makecell{45년} & \makecell{30년} \\
\hline
\end{tabular}\par}

4) 본 건물의 감가수정은 만년감가를 기준으로 함. \par
5) 창고 전체와 점포 중 20m\textsuperscript{2}가 도시계획도로에 저촉됨. \par
6) 본 건물 중 점포 잔여부분 보수비는 건축법상 요구되는 시설개선비 2,000,000원을 포함하여 총 5,000,000원임(담보감정시 보수비는 고려하지 아니함). \par
<자료 5> 기타참고사항 \par
1) 물음(2) 작성시 감정평가서 필수적 기재사항 중 평가액의 산출근거 및 그 결정에 관한 의견서를 목적별로 구체적으로 작성하시오(단, 물음(1)에서 작성된 가격산출근거는 생략). \par
2) 기호1 토지의 지상에 소재하는 이대한 소유의 건물과 박조선 소유의 건물은 합법적인 건축물로서 소유자별로 각각 점유하고 있음. \par
3) 가격시점에 대해서는 별도로 제시받지 아니함. \par
4) 보상평가를 위한 사업개요 \par
① 사업의 종류 : 乙지구 택지개발사업 \par
② 택지개발예정지구 공람ㆍ공고일 : 2008. 12. 01 \par
③ 택지개발예정지구 지정ㆍ고시일 : 2009. 09. 15 \par
5) 보상가격산정을 위한 참고자료 \par
① 택지개발예정지구 공람ㆍ공고일 이후 당해 공익사업지구내 표준지공시지가의 평균변동률이 당해 시ㆍ군ㆍ구 전체의 표준지공시지가의 평균변동률보다 1.4배 높고 그 변동률 차이는 10\%정도임. \par
② 기호1, 2 토지 공히 일부만 택지개발예정지구에 편입되었으나, 잔여지 손실보상에 대해서는 고려하지 아니함. \par
6) 감정평가사인 김한국씨는 감정평가서를 최종적으로 2010. 09. 04에 작성ㆍ완료하여, 2010. 09. 05에 심사 후 발송하였음. \par
\end{longquestion}
